% Generated draft: requires mathematical and source review.
\documentclass[11pt]{article}
\usepackage[margin=25mm]{geometry}
\usepackage{fontspec}
\setmainfont{Latin Modern Roman}
\usepackage{amsmath,amssymb,mathtools,unicode-math}
\setmathfont{Latin Modern Math}
\usepackage{enumitem,needspace,xurl,hyperref}
\setlength{\parindent}{0pt}
\setlength{\parskip}{5pt}
\setlength{\emergencystretch}{3em}
\newcommand{\N}{\mathbb{N}}
\newcommand{\Z}{\mathbb{Z}}
\newcommand{\Q}{\mathbb{Q}}
\newcommand{\R}{\mathbb{R}}
\newcommand{\C}{\mathbb{C}}
\newcommand{\F}{\mathbb{F}}
\newcommand{\A}{\mathbb{A}}
\newcommand{\PP}{\mathbb{P}}
\newcommand{\E}{\mathbb{E}}
\newcommand{\eps}{\varepsilon}
\newcommand{\dd}{\,\mathrm{d}}
\begin{document}
\section*{UnsolvedMath ID: 30005128}
\textbf{Source problem code:} catalogue entry
\subsection{Definitions and mathematical statement}
Fix a prime \(p\) and work in the category of \(p\)-local spectra. These are standard conventions for interpreting the supplied question; the record does not explicitly specify the prime or localization conventions. For an integer \(k\geq 1\), let \(T(k)\) denote the telescope of a \(v_k\)-self-map of a finite \(p\)-local spectrum of chromatic type \(k\). Here a finite spectrum \(F\) has type \(k\) if \(K(i)_*F=0\) for \(0\leq i<k\) and \(K(k)_*F\neq 0\), where \(K(i)\) is Morava \(K\)-theory at \(p\), with \(K(0)=H\Q\). Put \(T(0)=H\Q\). For a spectrum \(E\), write \(L_E\) for Bousfield localization at \(E\). A spectrum \(X\) is \(T(k)\)-acyclic if \(T(k)\wedge X\simeq 0\), equivalently \(L_{T(k)}X\simeq 0\).

For \(r\geq 0\), interpret finite chromatic localization by
\[
L^f_r=L_{T(0)\vee T(1)\vee\cdots\vee T(r)}.
\]
A spectrum \(X\) is \(L^f_r\)-local if its localization map \(X\longrightarrow L^f_rX\) is an equivalence. Let \(R\) be an algebra in spectra, as in the supplied excerpt, and let \(K(R)\) denote its algebraic \(K\)-theory spectrum. The record does not specify the precise multiplicative structure required of \(R\), or whether connective or nonconnective algebraic \(K\)-theory is intended.

\textbf{Extracted mathematical question.} In the standard indexing range \(n\geq 1\), does every \(L^f_{n-1}\)-local algebra in spectra \(R\), in the intended class of algebras, satisfy
\[
L_{T(m+1)}K(R)\simeq 0
\qquad\text{for every integer }m\geq n?
\]
Equivalently, is the implication
\[
R\xrightarrow{\simeq}L^f_{n-1}R
\quad\Longrightarrow\quad
\bigl(\forall m\in\Z,\ m\geq n:\ T(m+1)\wedge K(R)\simeq 0\bigr)
\]
valid for every integer \(n\geq 1\) and every such \(R\)? This displays the question under explicit conventional interpretations. The source does not state a range for \(n\); extending it to \(n=0\) would require a convention for \(L^f_{-1}\), which is not supplied. The intended class of algebras must also be confirmed before this is a fully specified universal statement.

\textbf{Broader source context and unresolved qualifications.} The excerpt distinguishes a qualitative chromatic Lichtenbaum--Quillen assertion from quantitative refinements. At height one, its qualitative assertion is that the fibre of
\[
K(R)_{(p)}\longrightarrow L^f_1K(R)_{(p)}
\]
is bounded above: for that fibre \(F\), there exists an integer \(B\) such that \(\pi_iF=0\) for every \(i>B\). Here the extracted notation \(K(R)(p)\) has been read as \(p\)-localization; the record does not explicitly confirm this reading or supply the hypotheses on \(R\). The quantitative refinement would specify a bound and might also describe \(L_{T(1)}K(R)\). No precise bound, proposed description, or complete higher-height formulation is supplied.

The excerpt also proposes two expressions deserving the name height:
\[
\inf\{j\in\Z:j\geq -1,\ L_{T(j+1)}R\simeq 0\},
\qquad
\sup\{j\in\Z:j\geq -1,\ L_{T(j)}R\not\simeq 0\}.
\]
These are transcribed as source expressions, not adopted as complete definitions: the second uses \(T(-1)\), which the record does not define, and neither expression comes with conventions for empty or unbounded indexing sets. The excerpt claims a general chromatic height-shift result and a stronger lower equivalence, but its referenced diagram and the exact equivalence are missing. Consequently, the extracted question must not automatically be treated as an unresolved conjecture.

The supplied background reports quantitative results for \(BP\langle 1\rangle\) and \(BP\langle 2\rangle\), and qualitative results for \(BP\langle n\rangle\) for general \(n\). These are source-reported claims about truncated Brown--Peterson spectra; the corresponding hypotheses and numbered references are not fully supplied. The record labels the broader problem partially solved and says that the full quantitative conjecture and all proposed refinements are not established. Those labels and assessments are supplied curation, not conclusions verified in this draft.
\subsection{Short English statement}
The extracted question asks whether finite chromatic locality of a ring spectrum \(R\) forces \(K(R)\) to be acyclic at every telescopic height \(m+1\) with \(m\geq n\). The displayed formulation uses explicit standard conventions, but the source leaves the class of ring spectra, some boundary conventions, and the broader quantitative conjecture unspecified. Its excerpt also reports a general height-shift result, so the supplied partially solved label cannot identify which displayed assertions remain open.
\subsection{Sources}
\begin{itemize}
\item Local UnsolvedMath registry, permanent ID \(30005128\), source problem code \texttt{catalogue entry}, title: Chromatic Lichtenbaum--Quillen and Redshift Conjectures. The supplied record belongs to the Oberwolfach Reports Open Problems collection and gives proposed year 2022. No individual proposer is supplied.
\item Supplied report citation: \emph{Algebraic K-Theory (2023)}, Oberwolfach Reports, pp.\ 1313--1376, DOI \texttt{10.4171/owr/2022/24}: \url{https://doi.org/10.4171/owr/2022/24}. The citation's year 2023 and the record's proposed year 2022 are retained as supplied; their relationship has not been verified.
\item Supplied literature source: Hahn--Wilson, redshift paper, \url{https://doi.org/10.4007/annals.2022.196.3.6}. The record does not supply its full bibliographic description or map it explicitly to the excerpt's numbered references.
\item Supplied literature source: Yang, \emph{Effective Redshift}, \url{https://arxiv.org/abs/2505.00344}.
\item Supplied literature source: Oberwolfach report, \url{https://doi.org/10.4171/owr/2022/24}.
\item Literature assessment dated exactly \texttt{2026-08-22}: supplied dated curation, not independently verified. It labels the problem partially solved and reports qualitative chromatic Lichtenbaum--Quillen and redshift results in broad generality, quantitative results in important low-height cases, and an incompletely established full quantitative conjecture and proposed refinements. The background explicitly requests independent verification of time-sensitive status and cited claims.
\item Attribution in the supplied excerpt: the term chromatic redshift is attributed to Ausoni and Rognes. The excerpt cites results using labels [2, 1] and [5], but the corresponding bibliography is absent. Its first-person assertion of a general height-shift result has no identified speaker in the supplied record.
\item The registry marks its statement \texttt{corrected\_verified} and reports that the official report was checked directly. This is a supplied verification claim; no source has been visited or independently checked for this draft.
\end{itemize}
\end{document}
